\documentclass[preview]{standalone}

\usepackage{amsmath}
\usepackage{amssymb}
\usepackage{stellar}
\usepackage{definitions}
\usepackage{tikz}
\usepackage{bettelini}

\begin{document}

\id{series-of-functions}
\genpage

\section{Series of functions}

\begin{snippetdefinition}{function-series-pointwise-convergence-definition}{Pointwise convergence of a series of functions}
    Let \(S\subseteq\realnumbers\) and let
    \(f_n\colon S\fromto\realnumbers\) for every \(n\in\naturalnumbers\).
    The \series \(\sum_{n=1}^{\infty}f_n\) \emph{converges pointwise} on
    \(S\) to \(s\colon S\fromto\realnumbers\) if the \partialsum[partial sums]
    \[
        S_n(x)=\sum_{k=1}^{n}f_k(x)
    \]
    converge pointwise to \(s\), that is,
    \[
        \forall x\in S,\qquad
        \lim_{n\to\infty}S_n(x)=s(x).
    \]
\end{snippetdefinition}

\begin{snippetproposition}{function-series-pointwise-cauchy-criterion}{Pointwise Cauchy criterion for series of functions}
    Let \(S\subseteq\realnumbers\) and let
    \(f_n\colon S\fromto\realnumbers\). The \series
    \(\sum_{n=1}^{\infty}f_n\) converges pointwise on \(S\) if and only if
    \[
        \forall x\in S,\ \forall\varepsilon>0,\ \exists N=N(x,\varepsilon)
        \in\naturalnumbers\ \suchthat
        \forall m>n\geq N,\qquad
        \left|\sum_{k=n+1}^{m}f_k(x)\right|<\varepsilon.
    \]
\end{snippetproposition}

\begin{snippetproof}{function-series-pointwise-cauchy-criterion-proof}{function-series-pointwise-cauchy-criterion}{Pointwise Cauchy criterion for series of functions}
    For a fixed \(x\in S\), the series converges at \(x\) precisely when the
    numerical sequence of partial sums \(\{S_n(x)\}\) is Cauchy. Since
    \[
        S_m(x)-S_n(x)=\sum_{k=n+1}^{m}f_k(x),
        \qquad m>n,
    \]
    the numerical Cauchy criterion gives the stated condition. Requiring it at
    every \(x\in S\) is exactly pointwise convergence.
\end{snippetproof}

\begin{snippet}{function-series-pointwise-cauchy-tail-remark}
    In the pointwise Cauchy criterion one may write \(m=n+h\), with
    \(h\in\naturalnumbers\). The index \(N\) may depend on the point \(x\),
    just as in pointwise convergence of sequences of functions.
\end{snippet}

\begin{snippetexample}{function-series-pointwise-discontinuous-sum-example}{A pointwise convergent series with discontinuous sum}
    On \(S=\realnumbers\), define
    \[
        f_n(x)=\frac{x^2}{(1+x^2)^n}.
    \]
    Every \(f_n\) is continuous. At \(x=0\), all terms vanish. If
    \(x\neq0\), then
    \begin{align*}
        \sum_{k=1}^{n}f_k(x)
        &=x^2\sum_{k=1}^{n}\left(\frac1{1+x^2}\right)^k,
    \end{align*}
    which is a geometric sum with ratio \(1/(1+x^2)\). Consequently,
    \[
        \sum_{n=1}^{\infty}f_n(x)
        =\begin{cases}
            0 & x=0,\\
            1 & x\neq0.
        \end{cases}
    \]
    Thus pointwise convergence does not preserve continuity. The same weak
    convergence generally does not justify interchanging differentiation or
    integration with summation.
\end{snippetexample}

\section{Uniform convergence}

\begin{snippetdefinition}{function-series-uniform-convergence-definition}{Uniform convergence of a series of functions}
    Let \(S\subseteq\realnumbers\) and let
    \(f_n\colon S\fromto\realnumbers\). The \series
    \(\sum_{n=1}^{\infty}f_n\) \emph{converges uniformly} on \(S\) to
    \(s\colon S\fromto\realnumbers\) if its partial sums
    \(S_n=\sum_{k=1}^{n}f_k\) converge uniformly to \(s\), equivalently if
    \[
        \forall\varepsilon>0,\ \exists N\in\naturalnumbers\ \suchthat
        \forall n>N,\ \forall x\in S,\qquad
        |S_n(x)-s(x)|<\varepsilon.
    \]
\end{snippetdefinition}

\begin{snippet}{function-series-uniform-convergence-supremum-formulation}
    Uniform convergence of \(\sum f_n\) to \(s\) is equivalently expressed by
    \[
        \sup_{x\in S}\left|\sum_{k=1}^{n}f_k(x)-s(x)\right|
        \longrightarrow0.
    \]
\end{snippet}

\begin{snippetproposition}{function-series-uniform-cauchy-criterion}{Uniform Cauchy criterion for series of functions}
    Let \(S\subseteq\realnumbers\) and let
    \(f_n\colon S\fromto\realnumbers\). The \series
    \(\sum_{n=1}^{\infty}f_n\) converges uniformly on \(S\) if and only if,
    for every \(\varepsilon>0\), there exists \(N\in\naturalnumbers\) such
    that
    \[
        \left|\sum_{k=n+1}^{m}f_k(x)\right|<\varepsilon,
        \qquad \forall m>n\geq N,\ \forall x\in S.
    \]
    Equivalently,
    \[
        \sup_{x\in S}\left|\sum_{k=n+1}^{m}f_k(x)\right|<\varepsilon,
        \qquad \forall m>n\geq N.
    \]
\end{snippetproposition}

\begin{snippetproof}{function-series-uniform-cauchy-criterion-proof}{function-series-uniform-cauchy-criterion}{Uniform Cauchy criterion for series of functions}
    Apply the uniform Cauchy criterion for sequences of functions to the
    sequence of partial sums \(S_n=\sum_{k=1}^{n}f_k\). For \(m>n\),
    \[
        S_m(x)-S_n(x)=\sum_{k=n+1}^{m}f_k(x),
    \]
    so its Cauchy condition is exactly the displayed tail estimate.
\end{snippetproof}

\end{document}
